% Clase 4 — Inferencia en LPO
% Lectura: AIMA Cap. 9
\section{Inferencia en LPO II}

\subsection{Encadenamiento hacia adelante (Forward Chaining)}

Aplica reglas de la forma $P_1 \land \cdots \land P_n \Rightarrow Q$ para derivar nuevos hechos hasta alcanzar la meta.
\begin{itemize}
    \item Completo para bases de conocimiento de \textbf{cláusulas de Horn}.
    \item Dirigido por los datos.
\end{itemize}

\subsection{Encadenamiento hacia atrás (Backward Chaining)}

Parte de la meta y trabaja hacia atrás buscando hechos o submetas.
\begin{itemize}
    \item Completo para cláusulas de Horn.
    \item Base de Prolog.
    \item Puede ciclar si no se controla.
\end{itemize}

\subsection{Resolución en LPO}

Generaliza la resolución proposicional usando unificación:

$$\frac{\ell_1 \lor \cdots \lor \ell_k \quad \neg m_1 \lor \cdots \lor \neg m_j}{\text{SUBST}(\theta, \ell_1 \lor \cdots \lor m_j)}$$

donde $\theta = \text{UNIFY}(\ell_i, m_k)$.

\begin{itemize}
    \item \textbf{Refutación por resolución}: completa para LPO (Teorema de Completitud de Gödel).
\end{itemize}

% Agrega aquí ejemplos de encadenamiento y resolución

